\chapter{Introduction}
\section{Physical Quantities}

Applications in many industry sectors rely on quantity values with units of measure, for sensor observations, actuation, design calculations and simulations, quantitative general knowledge, and clinical data. A temperature reading, for instance, is merely a bare number until it is paired with a specific unit such as degrees Celsius or Kelvin. The representation of such values is therefore critical: a numeric literal alone provides insufficient context, and associating each value correctly with its unit is a precondition for interoperable computation. To achieve this, the value and its unit must be encoded together, and any system that operates on them must respect the dimensional structure of the unit space. A precise foundation is therefore needed to define what quantities are, how they are classified, and how they combine.

The International Vocabulary of Metrology (VIM, JCGM 200:2012) makes the notion precise \cite{jcgm2012international}. It defines a quantity as a property of a phenomenon, body, or substance whose magnitude can be expressed as a number and a reference, where a reference can be a measurement unit. The definition separates two ideas that everyday usage tends to blur. A quantity kind, such as length or temperature, is general and qualitative. A quantity value is specific: the height of one particular object, stated as a number multiplied by a unit. Dimensional analysis governs how quantity values combine. Two quantities can be compared or aggregated only if they share the same dimension. Two measurements of the same dimension expressed in different units denote the same quantity kind. A system that treats them distinct is dimensionally wrong. A system that operates across them without conversion is arithmetically wrong. Representing physical quantities correctly therefore means carrying magnitude and unit together and respecting the dimensional structure of the unit at every step.


In software, this discipline is enforced by unit-aware libraries. Pint in Python and Indriya in Java are representative examples, and comparable libraries exist in most programming languages. In these libraries a quantity is a magnitude paired with a unit. An operation between incommensurable quantities raises an error at runtime, while compatible quantities are converted automatically when combined. Their scope is computation inside a single application. What these libraries do not address is how quantities are stored, serialized, and queried across distributed systems. Hall and Kuster identify representational difficulties with quantities and units in digital systems in general and propose a central metrological register to address them \cite{hall2022representing}.


On the Semantic Web the established answer is ontology-based. \ac{qudt} is an ontology, which grew out of a \ac{nasa} initiative, organises units within a hierarchy of quantity kinds and dimensions and records conversion factors. \ac{om} is also an ontology which provides a comparable unit vocabulary with finer-grained quantity kind hierarchies. \ac{sosa} and \ac{ssn}, published as a joint \ac{w3c} and \ac{ogc} standard, model the observation itself and delegate the encoding of value and unit to \ac{qudt} or \ac{om}. Keil and Schindler \cite{keil2018comparison} evaluated eight major unit ontologies and found coverage gaps and correctness issues across all of them. 

\section{The Problems}
Consider a dataset that records distances from several sources. A construction survey reports in inches, a road network in miles, and a scientific instrument in metres. A \ac{sparql} query asking which measurements exceed one kilometre ought to be straightforward. It is not. The query engine sees bare numbers next to unit IRIs and has no mechanism for treating them as commensurable quantities. Four problems lie behind this failure. Section 3 examines each in detail; this section states them in outline.

The first and most serious is computational inaccessibility. Ontologies such as \ac{qudt} and \ac{om} do encode dimensional structure and conversion factors, but that information never reaches the \ac{sparql} expression evaluator. Comparison, ordering, and aggregation operate on raw magnitudes, so queries across mixed units return physically incorrect results without warning. Lefrançois and Zimmermann \cite{lefranccois2018unified} confirm that, they are not aware of any existing \ac{sparql} engine provides unit-aware evaluation for \ac{qudt} or \ac{om} 
quantities. The only workaround is to make conversion factors explicit inside every query. 
At the time of that work, no such engine existed, and to the best of our knowledge 
that remains the case today.

The second is bounded coverage. An ontology enumerates its units one by one in a curated graph, yet the space of valid unit expressions has no natural boundary. Derived units multiply combinatorially: every pairing of a length unit with a time unit is a distinct speed unit, and the same holds for pressure, energy, and power. No enumeration can be complete.

The third is interoperability. Independently developed ontologies assign different IRIs to the same physical unit, and their vocabularies overlap only partially \cite{keil2018comparison}. Combining datasets from different sources typically requires cross-ontology mappings, which must be built manually and may not exist for a given pair of ontologies.

The fourth is verbosity. Encoding one quantity value costs several triples. In a dataset with thousands of observations the overhead grows linearly with the data while adding no expressiveness.

\section{cdt:ucum as a Datatype-Level Solution}
The \ac{lindt} Custom Datatypes specification approaches these four problems from a different level: not the ontology, but the \ac{rdf} datatype. It defines cdt:ucum, a datatype that encodes a complete quantity value as a single typed literal such as \texttt{"1.5 km"}\^{}\^{}\texttt{cdt:ucum}. The unit is part of the literal's lexical form, written in the \ac{ucum} grammar. A single triple suffices; no blank node is required. The specification defines operator semantics for equality, comparison, and arithmetic, exposed through native \ac{sparql} syntax, and its value space recognises dimensional equivalence, so two literals denoting the same physical quantity are value-equal regardless of the unit used to express them. Because \ac{ucum} is a grammar rather than a list, any unit the grammar can generate is valid. Coverage is unlimited by construction and independent of any curated vocabulary. None of this is achievable with ontology-based representations alone. Section 3 presents the specification and its formal datatype model in full.


\section{Prior Work and the Gap This Thesis Addresses:} Lefrançois and Zimmermann \cite{lefranccois2018unified} presented cdt:ucum at ESWC 2018 together with an implementation called jena-ucum. That implementation is a fork of the Apache Jena repository built against version 3, rather than a standalone extension. Using it means running the forked distribution instead of adding a dependency to a standard Jena installation. No other implementation exists for any other major \ac{rdf} library, and no general procedure for producing one has been documented.



\textbf{Research Question:} This thesis addresses the following research question: 
How can the implementation of cdt:ucum in an \ac{rdf} library be generalized into a procedure that is applicable and reproducible across different \ac{rdf} libraries?

\textbf{Contributions:} This thesis makes three contributions. The first is a five-step methodology for implementing cdt:ucum in any \ac{rdf} library. The second is four working implementations spanning Java (Apache Jena), TypeScript (rdflib.js), Python (RDFLib), and C\# (dotNetRDF), each accompanied by a test suite that reflects what that library and \ac{ucum} backend combination can and cannot achieve. The third is an empirical analysis of how \ac{rdf} library architecture and \ac{ucum} backend capabilities shape implementation strategy, documenting the limitations encountered in each case study and the patterns that recur across them.


\textbf{Outline of Remaining Sections:}
Section 2 covers background on \ac{rdf}, typed literals, and \ac{sparql}. Section 3 surveys the representation of physical quantities on the Semantic Web, covering ontology-based approaches, the \ac{ucum} code system, and the CDT specification. Section 4 reviews prior custom datatype implementations and the candidate \ac{rdf} and \ac{ucum} libraries. Section 5 presents the five-step methodology. Section 6 reports the four implementations and their findings. Section 7 concludes with a summary of contributions, stated limitations, and directions for future work.